\documentclass{article}
\usepackage{graphicx}
\input{preamble}

\title{Zadanie 4 . \\ 3D Gaussian Splatting}
\date{}

\begin{document}

\maketitle

\section*{Znenie zadania}

Cieľom zadania je oboznámiť sa s modernou metódou 3D rekonštrukcie scény pomocou
\textbf{3D Gaussian Splatting} (3DGS).
Študent prejde celým pipeline-om od fotografií až po interaktívny 3D viewer,
najprv na referenčnom datasete a potom na vlastnoručne nasnímanej scéne.
\\[4pt]
Zadanie je dokopy za \textbf{25~b.} a pozostáva z nasledujúcich úloh:

\begin{enumerate}

% ── ÚLOHA 1 ──────────────────────────────────────────────────────────────────
\item \textbf{[5~b.]} Spustenie pipeline-u na referenčnom datasete.\\
Stiahnite niektorý z verejne dostupných datasetov (napr.\ \textit{Tanks and
Temples – truck}, \textit{garden}, alebo iný kompatibilný dataset).
Z datasetu použite \textbf{iba fotografie} (adresár \texttt{input/}) –
\emph{žiadne predpočítané dáta}.
Celý pipeline spustite sami od začiatku podľa návodu \texttt{NAVOD.md}:       

\begin{enumerate}
    \item \texttt{./run.sh} – spustenie Docker kontajnera.
    \item \texttt{colmap gui} – spustenie COLMAP GUI a postupné vykonanie:
        \begin{itemize}
            \item Feature extraction (model: \texttt{SIMPLE\_PINHOLE})
            \item Feature matching (Exhaustive)
            \item Reconstruction (Sparse)
            \item Export modelu do \texttt{distorted/sparse/0/}
        \end{itemize}
    \item \texttt{python3 convert.py -s data/<dataset>} – undistortion fotografií.
    \item \texttt{python3 train.py -s data/<dataset>} – trénovanie 3DGS modelu
          (30\,000 iterácií, checkpointy po 7\,000 a 30\,000).
    \item \texttt{SIBR\_gaussianViewer\_app -m output/<uuid>} – vizualizácia výsledku.
\end{enumerate}

Zdokumentujte postup screenshotmi z každého kroku (COLMAP GUI, progress trénovania,
viewer). Uveďte dosiahnutú metriku \textbf{PSNR} po 7\,000 a 30\,000 iteráciách.

% ── ÚLOHA 2 ──────────────────────────────────────────────────────────────────
\item \textbf{[8~b.]} Snímanie vlastnej scény kamerou Ximea.\\
Vyberte a nasnímate vlastnú scénu priemyselnou kamerou Ximea (vysoké rozlíšenie).
Fotografie \textbf{downscalujte} pred spustením COLMAP
(odporúčame dlhšiu stranu na max.\ \textbf{600–1200~px}) – kamera má príliš
vysoké rozlíšenie na priame spracovanie a zároveň COLMAP aj 3DGS by bežali
neúnosne dlho.

Spustite celý pipeline rovnako ako v úlohe 1.
Uveďte dosiahnutú metriku PSNR a vizuálne porovnajte výsledok
s referenčným datasetom.

% ── ÚLOHA 3 ──────────────────────────────────────────────────────────────────
\item \textbf{[4~b.]} Vyhodnotenie výsledkov.\\
Porovnajte výsledky oboch scén (referenčný dataset vs.\ vlastná scéna):
\begin{itemize}
    \item tabuľka PSNR po 7\,000 a 30\,000 iteráciách,
    \item vizuálne porovnanie renderov z rovnakých uhlov pohľadu,
    \item komentár – čo ovplyvnilo kvalitu vlastnej scény
          (počet fotografií, osvetlenie, typ objektu, \ldots).
\end{itemize}

% ── ÚLOHA 4 ──────────────────────────────────────────────────────────────────
\item \textbf{[4~b.]} Teoretické vysvetlenie metódy.\\
V dokumentácii vysvetlite princípy (vlastnými slovami, nie copy-paste):
\begin{itemize}
    \item Čo sú 3D Gaussiány a aké majú parametre ($\mu, \Sigma, \alpha$, farba).
    \item Čo robí COLMAP (SfM) a prečo je potrebný.
    \item Ako funguje splatting – projekcia do 2D a alpha blending.
    \item Čo optimalizuje loss funkcia počas trénovania.
\end{itemize}

% ── ÚLOHA 5 ──────────────────────────────────────────────────────────────────
\item \textbf{[4~b.]} Dokumentácia.\\
Odovzdajte správu (PDF) obsahujúcu:
\begin{itemize}
    \item popis každého kroku pipeline-u s príslušnými screenshotmi,
    \item popis procesu snímania vlastnej scény (čo, kde, ako),
    \item grafické porovnanie výsledkov,
    \item záver s odporúčaniami – čo by ste urobili inak pre lepšie výsledky.
\end{itemize}

\end{enumerate}

\textbf{Termín odovzdania:} 12.\ cvičenie semestra

\newpage

% =============================================================================
\section*{Poznámky, tipy a pomôcky}

\subsection*{Postup – dodržujte NAVOD.md}

Všetky konkrétne príkazy, cesty a parametre sú detailne popísané v súbore
\texttt{NAVOD.md} v repozitári.
Pri akomkoľvek probléme si najprv prečítajte sekciu \emph{Riešenie problémov}
v návode.

% ── Dataset ──────────────────────────────────────────────────────────────────
\subsection*{Kde nájsť referenčný dataset}

\begin{itemize}
    \item \textbf{Tanks and Temples} (truck, train):
          \href{https://www.tanksandtemples.org/download/}{tanksandtemples.org/download}
    \item \textbf{Mip-NeRF 360} (garden, room, counter, \ldots):
          \href{https://jonbarron.info/mipnerf360/}{jonbarron.info/mipnerf360}
    \item Z datasetu použite \textbf{iba priečinok s fotografiami} –
          akékoľvek predpočítané sparse/dense dáta odstráňte.
\end{itemize}

% ── Downscale ─────────────────────────────────────────────────────────────────
\subsection*{Prečo a ako downscalovať fotografie}

Kamera Ximea produkuje snímky s rozlíšením až 12–20~Mpx.
COLMAP extrahuje príznaky na úrovni pixelov a pri vysokom rozlíšení:
\begin{itemize}
    \item feature extraction trvá rádovo hodiny,
    \item pamäť GPU pri trénovaní je výrazne vyššia.
\end{itemize}
\textbf{Odporúčané rozlíšenie:} dlhšia strana \textbf{1600–2000~px}
(zodpovedá cca 2–3~Mpx).
Zachovajte pomer strán, nestrihajte.
Downscale spustite \emph{pred} spustením COLMAP.

% ── Tipy na snímanie ─────────────────────────────────────────────────────────
\subsection*{Tipy pre kvalitné snímanie scény}

\begin{itemize}
    \item \textbf{Počet fotografií:} minimálne \textbf{60–120 snímok}.
          Viac fotografií = lepšia rekonštrukcia, ale dlhší čas COLMAP a trénovania.
    \item \textbf{Pokrytie:} pohybujte sa \emph{okolo celého objektu} v kruhu,
          ideálne na viacerých výškach (nízko, v strede, trochu zhora).
          Každý bod scény musí byť viditeľný z \textbf{aspoň 3 rôznych uhlov}.
    \item \textbf{Prekrytie snímok:} po každom kroku pootočte kameru max.\
          o 10–15° – susedné snímky musia mať aspoň 60\,\% spoločnej plochy.
    \item \textbf{Statická scéna:} počas celého snímania sa nesmie nič hýbať
          (žiadni ľudia v zábere, žiadny vietor, pohyblivé predmety odložte).
    \item \textbf{Osvetlenie:} rovnomerné difúzne svetlo (oblačný deň alebo
          interiér s viacerými zdrojmi).
          Vyhnite sa ostrým tieňom a priamemu slnku – tiene sa menia
          pri pohybe kamery a mätú rekonštrukciu.
    \item \textbf{Typ objektu:} ideálny objekt má \emph{textúrovaný, matný}
          povrch (napr.\ socha, nábytok, kôš, rastlina).
          \textbf{Problematické} sú:
          \begin{itemize}
              \item lesklé / zrkadlové povrchy (kov, sklo),
              \item homogénne plochy bez textúry (biela stena, jednofarebný stôl),
              \item priehľadné objekty (sklo, plast).
          \end{itemize}
    \item \textbf{Ohnisková vzdialenosť:} používajte fixné nastavenia kamery
          (expozícia, ISO, zaostenie) počas celého snímania.
          Automatické zaostenie medzi snímkami spôsobuje nekonzistenciu.
    \item \textbf{Rozostrenie a pohyb:} fotografie musia byť ostré.
          Pri slabom osvetlení zvýšte ISO namiesto predlžovania expozičného času.
\end{itemize}

% ── Výber scény ──────────────────────────────────────────────────────────────
\subsection*{Odporúčané typy scén}

\begin{itemize}
    \item Stolový objekt (soška, rastlina, hrnček) na neutrálnom pozadí –
          jednoduché na obísť, krátky čas snímania.
    \item Kútik miestnosti s nábytkom – väčší objem, zaujímavejší výsledok.
    \item Schodisko alebo chodba – dobrá textúra, ale vyžaduje viac snímok.
\end{itemize}

% ── Metriky ──────────────────────────────────────────────────────────────────
\subsection*{Interpretácia metriky PSNR}

\begin{center}
\begin{tabular}{ll}
\toprule
PSNR & Kvalita \\
\midrule
$< 22$~dB & slabá rekonštrukcia (málo snímok, zlé pokrytie) \\
$22$–$26$~dB & prijateľná (checkpoint 7\,000 it.) \\
$26$–$29$~dB & dobrá (checkpoint 30\,000 it.) \\
$> 29$~dB & výborná \\
\bottomrule
\end{tabular}
\end{center}

% ── Užitočné odkazy ──────────────────────────────────────────────────────────
\subsection*{Užitočné odkazy}

\begin{itemize}
    \item Originálny paper: Kerbl et al., \textit{3D Gaussian Splatting for
          Real-Time Radiance Field Rendering}, SIGGRAPH 2023.\\
          \href{https://repo-sam.inria.fr/fungraph/3d-gaussian-splatting/}{repo-sam.inria.fr/fungraph/3d-gaussian-splatting}
    \item Vizuálny tutoriál:
          \href{https://learnopencv.com/3d-gaussian-splatting/}{learnopencv.com/3d-gaussian-splatting}
    \item COLMAP dokumentácia:
          \href{https://colmap.github.io/}{colmap.github.io}
    \item Repozitár so zdrojovým kódom a návodom: \texttt{NAVOD.md}
\end{itemize}

\end{document}
